\chapter*{ЗАКЛЮЧЕНИЕ}
\addcontentsline{toc}{chapter}{ЗАКЛЮЧЕНИЕ}

В рамках курсовой работы был разработан компилятор языка B на языке C. Компилятор выполняет чтение исходного текста, лексический анализ, синтаксический разбор, построение абстрактного синтаксического дерева, ведение таблицы символов и генерацию ассемблерного кода x86-64 для операционной системы Linux.

В аналитической части была рассмотрена грамматика языка B и выделены его основные конструкции: внешние определения, функции, автоматические и внешние переменные, векторы, адресная арифметика, операторы \texttt{if}, \texttt{while}, \texttt{switch}, \texttt{goto}, \texttt{return}, выражения и литералы. Также были сформулированы требования к компилятору.

В конструкторской части была разработана структура транслятора. Компилятор разделен на модули лексического анализа, синтаксического анализа объявлений, операторов и выражений, работы с AST, таблицей символов, представлением машинных слов и генерацией кода. Такое разделение позволило изолировать машинно-зависимую часть генератора от общего обхода дерева.

В технологической части была описана реализация компилятора, приведены примеры программ на языке B, фрагмент сгенерированного ассемблерного кода и результаты тестирования. Тестовые программы проверяют рекурсию, циклы, сортировку, обход графов, глобальные и локальные векторы, указатели, тернарный оператор, \texttt{switch}, \texttt{goto}, внешние объявления, комментарии и escape-последовательности.

Таким образом, были решены следующие задачи:

\begin{itemize}
  \item изучена документация языка B;
  \item описана грамматика поддерживаемого языка;
  \item реализован лексический анализатор;
  \item реализован синтаксический анализатор объявлений, операторов и выражений;
  \item реализована таблица символов;
  \item реализовано построение AST;
  \item реализована генерация x86-64 assembly;
  \item обеспечена сборка исполняемого файла под Linux;
  \item подготовлены и описаны регрессионные тесты.
\end{itemize}

Цель курсовой работы достигнута: разработан компилятор языка B, транслирующий исходные программы в исполняемый машинный код для архитектуры x86-64 под операционную систему Linux.
